\section{Môi trường và công nghệ}

\subsection{Công nghệ sử dụng}

Bảng~\ref{tab:deployment-environment} liệt kê các công nghệ chính. Phiên bản
image hạ tầng được ghim trong tệp Docker Compose để môi trường có thể tái tạo.

\begin{table}[H]
\centering
\caption{Các thành phần phần mềm của hệ thống}
\label{tab:deployment-environment}
\begin{tabularx}{\textwidth}{L{3.6cm}Y}
\toprule
\textbf{Thành phần} & \textbf{Công nghệ và phiên bản}\\
\midrule
Ngôn ngữ backend & Python 3.12, lập trình bất đồng bộ \texttt{asyncio}\\
Backend/API & FastAPI 0.141, Uvicorn, Pydantic 2.13, pydantic-settings\\
Database & PostgreSQL 16 (image \texttt{postgres:16-alpine}), driver \texttt{asyncpg}\\
Cache telemetry & Redis 7 (\texttt{redis:7-alpine})\\
Message broker & Eclipse Mosquitto 2 (MQTT), client paho-mqtt qua gói \texttt{hopquangdo-messaging}\\
Model store & Hugging Face Hub (\texttt{huggingface\_hub})\\
Container registry & GitHub Container Registry (GHCR)\\
Container engine (Edge) & Docker, NVIDIA Container Runtime trên Jetson\\
AI Runtime & Ultralytics YOLO (PyTorch/ONNX/TensorRT), OpenCV\\
AI Assistant & LangGraph 1.x, LangChain 1.x, MCP Python SDK, Claude Sonnet 5 (mô hình chính), Claude Haiku 4.5 (mô hình phụ)\\
Giám sát & Prometheus 2.54, Alertmanager 0.27, Grafana 11.2, node-exporter 1.8, cAdvisor 0.49, DCGM exporter 3.3\\
Frontend & React 19, TypeScript 6, Vite 8, Tailwind CSS 4, React Router 7\\
Kiểm thử và CI/CD & pytest, pytest-asyncio, Ruff, oxlint, GitHub Actions\\
\bottomrule
\end{tabularx}
\end{table}

\subsection{Môi trường Cloud}

Cloud được triển khai trên một máy ảo AWS EC2 chạy Linux, tại thư mục
\path{/opt/edgeops}. Các thành phần được đóng gói thành bốn \emph{stack} Docker
Compose độc lập, dùng chung mạng Docker \texttt{edgeops-cloud}:

\begin{itemize}
  \item \textbf{cloud-infra:} PostgreSQL, Redis (có mật khẩu), Mosquitto. Các tệp
        migration SQL được gắn vào \path{/docker-entrypoint-initdb.d} nên schema
        được tạo tự động ở lần khởi động đầu tiên.
  \item \textbf{cloud:} image \texttt{cloud-api}, một tiến trình phục vụ REST,
        WebSocket, MCP và embedded worker trên cổng 8000.
  \item \textbf{assistant:} image \texttt{assistant} trên cổng 8100, gọi Cloud qua
        tên dịch vụ nội bộ \texttt{http://api:8000}.
  \item \textbf{web:} image \texttt{web-dashboard} (Nginx) trên cổng 80, phục vụ
        giao diện và reverse proxy tới API, WebSocket và Assistant.
\end{itemize}

Stack \texttt{cloud-monitoring} (Prometheus, Alertmanager, Grafana) được khởi chạy
riêng khi cần giám sát metric hạ tầng. Tất cả cấu hình nhạy cảm (mật khẩu, API
key, token Hugging Face, khóa LLM) được đọc từ tệp \texttt{.env}, không nằm trong
mã nguồn.

\subsection{Môi trường Edge}

Mỗi Edge Node cần Docker (và NVIDIA Container Runtime nếu dùng GPU). Stack Edge
chỉ chạy tĩnh một container \texttt{edge-agent}; container này được gắn Docker
socket của máy chủ để tạo các container AI Runtime anh em (sibling container).
Cấu hình node gồm \texttt{EDGE\_CLUSTER\_ID}, \texttt{EDGE\_NODE\_ID}, địa chỉ
MQTT và API key. Image Agent được build đa kiến trúc \texttt{linux/amd64} và
\texttt{linux/arm64}, chạy được trên máy x86 và Jetson; image runtime cho Jetson
cần base image L4T tương ứng phiên bản JetPack. Hai thiết bị mục tiêu là Jetson
Nano và Jetson Orin Nano (thông số ở Bảng~\ref{tab:jetson-hardware}); máy trạm x86
có GPU NVIDIA được dùng làm Edge Node trong quá trình phát triển.

\subsection{Mạng}

Edge chủ động kết nối ra Cloud qua MQTT (cổng công bố 11883) nên không cần mở
cổng vào trên thiết bị; điều này phù hợp với thiết bị nằm sau NAT. Tải model đi
qua HTTPS tới Hugging Face Hub; metric hạ tầng được đẩy qua HTTPS bằng
Prometheus remote write. Script \texttt{tunnel-cloud-infra.sh} tạo đường hầm SSH
tới PostgreSQL, Redis, MQTT trên Cloud để phát triển từ máy cá nhân mà không công
khai các cổng quản trị.

\subsection{CI/CD}

Quy trình tích hợp và triển khai liên tục dùng GitHub Actions
(Bảng~\ref{tab:ci-workflows}). Workflow tái sử dụng \texttt{\_build-deploy.yml}
build image, đẩy lên GHCR với tag là commit SHA, sao chép tệp Compose và \texttt{.env}
(lấy từ GitHub Secrets) lên EC2 qua SSH, chạy \texttt{docker compose pull} và
\texttt{up -d}, rồi gọi health check sau khi cập nhật.

\begin{table}[H]
\centering
\caption{Các workflow CI/CD}
\label{tab:ci-workflows}
\begin{tabularx}{\textwidth}{L{3.3cm}L{3.6cm}Y}
\toprule
\textbf{Workflow} & \textbf{Kích hoạt} & \textbf{Công việc}\\
\midrule
ci.yml & Mọi pull request, push lên main & Ruff, pytest cho backend; lint và build Dashboard\\
ci-edge.yml & Thay đổi mã Edge & Build image \texttt{edge-agent} (amd64 + arm64) và \texttt{edge-runtime}; push lên GHCR khi merge\\
deploy-cloud.yml & Thay đổi Cloud, migration & Build \texttt{cloud-api}, triển khai, kiểm tra \texttt{/health}\\
deploy-assistant.yml & Thay đổi Assistant & Build và triển khai \texttt{assistant}\\
deploy-web.yml & Thay đổi Dashboard & Build và triển khai \texttt{web-dashboard}\\
\bottomrule
\end{tabularx}
\end{table}
